\SourceLeaf{351}%
Hilfskugel, das diesem Dreieck als ganze Krümmung entspricht, zur gesamten Oberfläche der Hilfskugel. Allgemein ist der Überschuss der Winkel eines Vielecks mit \(n\) Seiten über \(2n-4\) rechte Winkel gleich der ganzen Krümmung des Vielecks, wenn seine Seiten kürzeste Linien sind.

Die allgemeinen Untersuchungen der Abhandlung werden an ihrem Schluss noch auf die Theorie der aus kürzesten Linien gebildeten Dreiecke angewandt. Davon führen wir hier nur einige Hauptsätze an. Seien \(a,b,c\) die Seiten eines solchen Dreiecks, aufgefasst als Größen erster Ordnung, \(A,B,C\) die gegenüberliegenden Winkel, \(\alpha,\beta,\gamma\) die Krümmungsmaße in den Eckpunkten und \(\sigma\) der Flächeninhalt des Dreiecks. Dann ist, bis auf Größen vierter Ordnung, \(\tfrac13(\alpha+\beta+\gamma)\sigma\) der Überschuss der Summe \(A+B+C\) über zwei rechte Winkel. Mit derselben Genauigkeit lauten ferner die Winkel eines ebenen geradlinigen Dreiecks mit den Seiten \(a,b,c\) der Reihe nach
\[
\begin{aligned}
&A-\tfrac1{12}(2\alpha+\beta+\gamma)\sigma\\
&B-\tfrac1{12}(\alpha+2\beta+\gamma)\sigma\\
&C-\tfrac1{12}(\alpha+\beta+2\gamma)\sigma
\end{aligned}
\]
Man sieht sofort, dass der letzte Lehrsatz den bekannten, zuerst von \textsc{Legendre} aufgestellten Satz verallgemeinert. Nach diesem erhält man bis auf Größen vierter Ordnung die Winkel des geradlinigen Dreiecks, indem man jeden Winkel des sphärischen Dreiecks um ein Drittel des sphärischen Exzesses vermindert. Auf einer nichtsphärischen Fläche muss man also ungleiche Winkelkorrekturen anbringen. Die Ungleichheit ist im Allgemeinen eine Größe dritter Ordnung. Weicht jedoch die ganze Fläche nur wenig von der Kugelgestalt ab, so enthält diese Ungleichheit zusätzlich einen Faktor von der Ordnung der Abweichung von der Kugelgestalt. Für die höhere Geodäsie ist es zweifellos wichtig, die Ungleichheiten dieser Korrekturen berechnen zu können und sich dadurch ganz davon zu überzeugen, dass sie für alle messbaren Dreiecke auf der Erdoberfläche völlig unmerklich sind. So ergeben sich zum Beispiel im größten Dreieck der vom Verfasser ausgeführten Triangulation, dessen größte Seite fast 15 geografische Meilen lang ist und dessen Winkelsumme zwei rechte Winkel um fast 15 Sekunden übersteigt, für die drei Reduktionen der Winkel auf die %
